% Sections of Chapter 2 (Background). The lead-in lives in
% content/background/background.tex.
\section{Temporal Abstraction and Automated Support to Guideline-Based Care}
\label{sec:rw-guidelines}

The work reviewed in this section forms one continuous line of research, developed largely at the BGU Medical Informatics Research Center, running from the representation of clinical concepts to the retrospective assessment of clinical practice. It is the source of both the abstraction layer and the adherence signal used in this thesis.

That line begins with \textbf{knowledge-based temporal abstraction} (KBTA) \citep{shahar1997framework}, the framework that turns time-stamped clinical data into named interval concepts, and with the \textbf{Temporal Mediator} \citep{shahar1997framework}, the reasoning engine that executes a TAK knowledge base against a patient record to derive them (Section~\ref{sec:prelim-kbta}). The Mediator is the point at which the rest of this chapter connects: the interval streams mined by KarmaLego, aligned by iDTW, and scored for guideline adherence are all its output. It is used in this thesis through our own reimplementation \citep{mediator-github}.

\textbf{Clinical guidelines} (GLs) are evidence-based protocols developed to standardize medical care, improve treatment outcomes, and reduce variability in clinical decisions. Their implementation in practice, however, is complicated: GLs are text-based and not easily accessible to clinicians who must adapt them to individual patients and apply them at the point of care, which leads to low adherence among physicians. Consequently, there is a need for clinical decision support systems (CDSSs) that automatically apply and assess GLs, and for acquiring those GLs in a machine-interpretable form that can be enforced algorithmically. Doing so requires formalizing two kinds of knowledge: the declarative knowledge above, which defines what the guideline's terms mean, and \emph{procedural} knowledge, which specifies what to do and when.

\textbf{Asbru} \citep{asbru-task-specific-intention-based-time-oriented-language, asgaard-project-task-specific-framework-application-critiquing-time-oriented-clinical-guidelines} is the computer-interpretable language in which that procedural half is conventionally expressed, representing sequential, concurrent, conditional, and repetitive actions as a hierarchy of execution plans. \textbf{DeGeL} \citep{supporting-episodic-application-clinical-guidelines}, the Digital Electronic Guidelines Library, is a web-based distributed architecture that facilitates the gradual conversion of guidelines from text into a formal Asbru representation, handled through a graphical tool, \emph{Gesher}, with which a knowledge engineer and a clinical expert construct the specification together.

\textbf{Picard} \citep{architecture-continuous-user-driven-data-driven-application-clinical-guidelines} brings the two halves together into a continuous-care guideline application framework, combining declarative knowledge executed by the Mediator with procedural knowledge in Asbru, both drawn from DeGeL, to generate real-time recommendations, adherence reports, and alerts for critical deviations from a patient's longitudinal record. It served as the core decision-support engine of the \textbf{MobiGuide} project \citep{mobiguide-guiding-clinicians-and-chronic-patients-anytime-anywhere}, which supported chronic disease management in Spain and Italy by integrating hospital records, mobile biosensors, and patient inputs. There, the engine improved clinicians' compliance with guidelines from 41\% in manual mode to 93\% --- direct evidence that adherence is malleable, and therefore worth measuring.

\textbf{e-Picard} \citep{supporting-episodic-application-clinical-guidelines} adapts Picard to non-continuous, episodic care, where treatment is sought sporadically rather than monitored in real time. It conducts retrospective quality assessment of past actions using temporal fuzzy logic, which required a third body of knowledge alongside the declarative and procedural ones. \textbf{Quality-Assessment (QA) knowledge} re-expresses the procedural content of a guideline declaratively, as the temporal pattern its correct execution would have left in the record --- a Quality-Assessment Temporal Pattern (QATP). Acquiring it is a knowledge-engineering task in its own right: each plan in the procedural specification must be classified by the kind of constraint it imposes, given a membership function that grades partial compliance rather than issuing a binary verdict, and placed in a weighted tree mirroring the guideline's plan hierarchy through which scores are aggregated (Section~\ref{sec:prelim-qa}). Because QATPs are written in TAK, the Mediator computes adherence scores in the same pass that produces ordinary abstractions.

Systems of this kind are typically evaluated on two criteria --- \emph{correctness}, the accuracy of the recommendations produced, and \emph{completeness}, their coverage of guideline-recommended actions --- assessed first through technical evaluation against controlled, simulated scenarios and then through silent evaluation on real-world data, where recommendations are withheld from the treating team and reviewed independently by experts \citep{multiple-scenario-assessment-guideline-based-decision-support-compliance}. The QA machinery itself has been applied to real institutional settings: a system for assessing long-term management at a geriatric hospital \citep{implementation-evaluation-long-term-geriatric-hospital-quality} showed that automated assessment recovers clinically meaningful variation in how protocols are followed over extended periods.

This thesis adopts the QATP formulation of adherence, but not the systems that introduced it. No recommendations are produced and no decision-support engine is deployed: scoring is performed offline, over records already collected, by our own implementation. Three further differences position the work. In \emph{representation}, rather than maintaining a procedural formalism alongside a declarative one, both are encoded in a single TAK-based schema evaluated by one engine; a further contribution is the knowledge base itself, which formalizes the current national inpatient diabetes guidelines together with their QA extension, and is transferable to any EHR within that clinical scope. In \emph{scale}, prior assessments have generally been confined to modest cohorts within a single institution, which cannot separate the practice of a particular ward from the effect of a guideline, nor observe how practice shifts when a guideline is introduced nationally; here adherence is measured across six hospitals spanning such an implementation. In \emph{purpose}, these systems are evaluated \emph{against} the guideline --- it is the yardstick, and its clinical value is inherited from the research that established it. This thesis encodes the same predetermined guideline, but treats its individual recommendations as hypotheses to be tested against what subsequently happens to the patient, and uses the resulting scores as features for outcome prediction. Adherence thus serves as an explanatory variable rather than only as the quantity reported: the question is not whether the guideline was applied correctly, but which of its recommendations demonstrably matter.

\section{Time Interval Mining}
\label{sec:rw-mining}

This section reviews the methods that turn a longitudinal EHR into a feature-vector representation of a patient's state up to a given point in time, and the studies that have used them for clinical prediction. The pipeline they compose supplies the symbolic baseline of this thesis; its machinery is defined in Section~\ref{sec:prelim-tirp}.

\textbf{KarmaLego} \citep{moskovitch2015karmalego} extends KBTA with time-interval mining, discovering frequent Time Interval Relation Patterns (TIRPs) across a population of abstracted records --- a period of high-dose insulin administration preceding two or more hypoglycemic episodes, for instance, wherever that arrangement recurs often enough. It has been shown that both temporal abstractions of raw data and combinations of them into more complex patterns serve as effective features for machine-learning tasks including clustering, classification, prediction, and the interactive analysis of large temporal datasets. The \textbf{Semantic Adjacency Criterion (SAC)} \citep{sac-time-intervals-mining} refines the mined output by discarding patterns whose constituent intervals are separated by intervening ones that break the relationship, reducing computational overhead substantially and improving the interpretability of what survives.

This machinery has been applied not only to forecasting outcomes but to predicting the next clinical \emph{action}, recommending what a clinician should do next from the patterns observed in the histories of similar patients. That framing carries an assumption worth making explicit, because this thesis shares it: predicting a next best action presupposes that an optimal action exists and is identifiable. Quality-assessment scoring rests on the same premise from the opposite direction --- rather than predicting the correct action, it measures how far the recorded care departed from one already specified by a guideline.

\textbf{Interval-Based Dynamic Time Warping (iDTW)} \citep{implementation-evaluation-multivariate-abstraction-based-interval-based-dtw} takes a different route through the same abstraction layer. Classical dynamic time warping measures similarity between time series of different lengths by finding an optimal alignment between them, but is limited to univariate sequences, assumes complete data, and does not consider clinical abstractions. iDTW extends it to interval-based, multivariate, abstracted records --- ``HIGH glucose for 2 days'' --- through an interval representation of states and trends and an interval-matching distance that accounts for both value and duration. The resulting distance matrix between patients was used by a KNN-like model to perform classification.

\textbf{Temporal Probabilistic Profiles (TPF)} \citep{sheetrit2019tpf} close the pipeline by summarizing the mined patterns into one fixed-length vector per patient, usable directly for classification by vector-space methods or as meta-features for more advanced models. In the same work, this methodology was used to predict sepsis complications in ICU patients. On MIMIC-III, a random forest trained on TPF vectors was compared with a GRU-D recurrent model \citep{che2018recurrent} trained on the original multivariate raw data together with demographics, and across various configurations the TPF-based approach paired with a traditional machine-learning model consistently outperformed the more complex sequential model. Despite a different dataset and prediction problem, this establishes the potential of the TPF methodology for representing the first $K$ days of an admission, and the KarmaLego--TPF pipeline accordingly serves as the symbolic baseline of this thesis.

Its limitation is representational rather than empirical. Summarizing a patient as a distribution over mined patterns discards the order and timing of the sequence that produced it, and the mining step commits in advance to which arrangements are worth counting. The baseline used here accordingly follows the TPF construction of the per-patient distribution vector, while the classifier trained on it is our own. INTERVenE takes a separate route from the same Mediator output: it performs no pattern mining, consuming the interval sequence directly, so that ordering and duration reach the model intact, addressing this limitation directly.

\section{Sequential Models in Medical Outcome Prediction}
\label{sec:rw-sequential}

Sequential models, particularly Transformers \citep{attention-is-all-you-need}, have emerged as a powerful approach for outcome prediction by capturing dependencies across temporal sequences. The attention mechanism addresses the limitations of recurrent models in handling long-range dependencies and is capable of capturing \emph{implicit} patterns in temporal data, unlike KarmaLego-based methods, which capture explicit ones. The studies reviewed below applied such models to raw EHR data for various predictive tasks, and provided insight into preprocessing strategies, embeddings, and positional encodings relevant to structuring the experiments in this research.

Before attention, recurrent architectures dominated. \textbf{GRU-D} \citep{che2018recurrent} handles irregular sampling by augmenting a gated recurrent unit with learned exponential decays over the time elapsed since a variable was last observed, pulling a stale value toward its empirical mean so that missingness and staleness are modeled rather than imputed away. It remains the canonical recurrent approach to irregular clinical series and serves as a baseline in this thesis. A further study \citep{treatment-prediction-icu-partitioned-sequential-deep-time-series-analysis} introduced a framework for ICU treatment prediction using a bidirectional LSTM over multivariate time series and static features from MIMIC-IV, with partitioned prediction windows --- a 12-hour horizon split into three 4-hour sub-windows, predicted from the previous 24 hours --- and a two-step model that first predicted treatment necessity and then estimated dosage for hypoglycemia, hypokalemia, and hypotension. The study demonstrated the effectiveness of temporal neural networks at mimicking ICU decision-making, but operates on raw measurements throughout, using neither abstraction nor any representation of how well the recorded care matched what was indicated.

Attention-based models subsequently removed the need for imputation altogether. \textbf{STraTS} \citep{tipirneni2022strats} abandons the fixed time grid: each observation is represented as a \emph{(variable, value, time)} triplet, in which the variable is embedded from a lookup, and the continuous value and timestamp are each embedded by a learned continuous encoding rather than discretized. The resulting unordered set of triplet embeddings is processed by a Transformer and fused into a patient representation by attention, so that sparsity and irregularity require no resampling, no alignment to a common grid, and no imputation of unobserved variables --- the difficulty GRU-D absorbs through decay is instead removed from the representation. The architecture is further trained with a self-supervised forecasting objective over unlabeled admissions before supervised fine-tuning, exploiting the large pool of admissions that carry no outcome label. Two configurations accordingly serve as baselines in this thesis: \textbf{STraTS (pre+ft)}, forecasting pretraining followed by supervised fine-tuning, and \textbf{ss-STraTS}, the same architecture trained supervised-only from scratch, which isolates how much the pretraining stage contributes on this cohort.

STraTS is the closest published counterpart to this thesis in modeling philosophy: both replace grid-aligned matrices with a set of semantically typed tokens carrying their own timestamps, and both leave the timing continuous rather than binning it. They differ in what a token \emph{is} --- a single raw measurement against an abstracted interval --- which is precisely the comparison the evaluation is designed to isolate. Both STraTS configurations, alongside GRU-D, therefore form the neural baselines against which INTERVenE is measured.

A parallel family operates at the visit level rather than the measurement level. \textbf{BEHRT} \citep{li2020behrt} treats diagnosis codes as tokens and visits as positions, while \textbf{HEART} \citep{huang2024heart} models heterogeneous relations among clinical entities explicitly. \textbf{TransformEHR} \citep{yang2023transformehr} presents a generative encoder--decoder architecture for predicting future disease outcomes from longitudinal EHR data, processing multi-level embeddings of ICD-10 codes, visit-level information, and demographics, and encoding time gaps between visits as sinusoidal embeddings of visit dates. Its cross-attention connects historical patient data to future disease predictions, and it is pretrained to recover all codes of a future visit from past visits --- a strategy that outperformed the conventional code-level masking of BERT-style models. The study highlights challenges similar to ours, particularly in handling multi-dimensional sequential EHR data and determining how best to encode time and integrate visit-level information into the architecture.

\textbf{ETHOS} \citep{zero-shot-health-trajectory-prediction-transformer} introduces a decoder-only transformer for zero-shot prediction of Patient Health Timelines. It treats health events --- admissions, laboratory tests, medications --- as tokens, embeds static patient data at the start of each sequence, and encodes temporal gaps using explicit quantile-based interval tokens rather than learned encodings of the exact $\Delta t$, enabling it to operate directly on irregular, noisy data. ETHOS achieves strong performance across ICU mortality, hospital mortality, length-of-stay estimation, and DRG classification without task-specific fine-tuning, underscoring the feasibility of a single trained model performing a broad range of patient-trajectory predictions.

Neither TransformEHR nor ETHOS incorporates temporal abstraction, pattern mining, or adherence scoring, and both quantize elapsed time --- into visit-date embeddings and into interval-token bins respectively --- rather than representing it continuously. Both choices are revisited in this thesis: Time2Vec \citep{kazemi2019time2vec} supplies a continuous encoding of $\Delta t$, and the token itself carries the abstraction these models leave to be inferred.

The property shared by all of these models is the granularity of their tokens: either point-in-time raw values or heavily aggregated visit-level codes. Neither carries the duration, trend, or clinical context of a state --- a glucose value of $260$ says nothing about whether it is the fourth such value in a rising sequence or an isolated excursion. Such structure can in principle be recovered from data, but only implicitly, and what is recovered cannot be inspected. Building on these insights, INTERVenE --- the transformer family developed in this thesis --- operates on abstracted intervals as tokens, so that duration, trend, and context are present in the input by construction. What that representation contributes separates into two questions the evaluation addresses independently: what the persistence of a symbol over time contributes to prediction, and what its naming contributes to interpretability.

\section{Explainable AI: Explaining a Model's Decision}
\label{sec:rw-xai}

Clinical trust in a predictive model requires more than accuracy. Post-hoc attribution methods --- SHAP \citep{unified-approach-interpreting-model-predictions} and LIME \citep{why-should-i-trust-you-explaining-predictions-any-classifier}, whose mechanics are described in Section~\ref{sec:prelim-xai} --- assign importance to individual features for a specific prediction, and can be applied to sequential architectures at the token or event level, quantifying how particular time steps influence the prediction of future events. This token-level attribution is critical for medical applications, where the interpretability requirement extends beyond knowing \emph{which} variables matter to understanding \emph{when} and \emph{in what context} they matter.

An instructive example of such methods in medical informatics is a multimodal model \citep{feature-importance-explain-multimodal-prediction-models-clinical-use-case} for predicting 30-day postoperative mortality in elderly hip-fracture patients, which combined pre-operative static data, radiographs, and per-operative vital signs and medications. Shapley values provided both global explanations --- showing that pre-operative static data accounted for roughly 70\% of the predictive power --- and individualized local explanations. Their use differs from ours in purpose: they identify important predictors of an outcome, whereas we aim to use such insights on temporal patterns to inform the guidelines themselves, though the underlying goal of producing actionable insight is shared.

A complementary line of work, developed within mechanistic interpretability research on large language models, looks inside the network rather than around it. Attention maps show how information is routed between tokens, but they do not decompose the semantic content that the MLP blocks compute over it, and individual hidden units are seldom interpretable in isolation. \textbf{Sparse autoencoders} \citep{rajamanoharan2024jumping} address this by learning an over-complete, sparsely activating dictionary over a layer's activations, resolving them into directions each active for few inputs and correspondingly easier to name. \textbf{Sparse transcoders} \citep{dunefsky2024transcoders} change what is learned: rather than reconstructing activations in place, they reproduce the input--output \emph{mapping} of a component, so the dictionary characterizes the computation rather than only its output and its features can be followed through the network into circuits. Both constructions are given in Section~\ref{sec:prelim-xai}.

The consequence that matters for a clinical setting is what this basis permits. Because the dictionary describes the network's own computation, an attribution built on it is computed against the deployed model rather than against a surrogate fitted around its outputs --- the standing objection to SHAP and LIME, whose explanations describe an approximation of the model rather than the model itself.

This machinery has seen little application to clinical prediction, and the reason is instructive: dictionary learning yields interpretable \emph{units}, but whether a unit means anything clinically depends entirely on what the network was fed. A dictionary learned over raw measurement embeddings resolves to directions that a clinician still has to interpret.

Both approaches therefore rest on the same precondition, and it is the one this thesis supplies. When every token is a named abstraction, an explanation resolves to a clinical statement by construction rather than by interpretation after the fact; and when guideline-adherence scores form part of that input, the explanations reach past patient state into the treatment process, identifying which recommendations drive predicted risk. Our goal is to use these insights to identify which aspects of the protocol are most critical to treatment outcomes, and ultimately to refine the protocol accordingly.
